\documentclass[14pt,a4paper]{extarticle}
\usepackage{fontspec}
\usepackage{polyglossia}
\setmainlanguage{russian}
\setmainfont{Times New Roman}
\setmonofont{Courier New}
\usepackage{setspace}
\onehalfspacing
\usepackage{geometry}
\geometry{left=30mm, right=15mm, top=20mm, bottom=20mm}
\usepackage{indentfirst}
\setlength{\parindent}{1.25cm}
\setlength{\parskip}{0pt}
\usepackage{xurl}
\urlstyle{same}
\usepackage{hyperref}
\usepackage{booktabs}
\usepackage{ragged2e}
\usepackage{graphicx}
\usepackage{float}
\usepackage{pdflscape}
\usepackage{listings}
\usepackage{xcolor}
\usepackage{array}
\usepackage{amsmath}
\usepackage{wasysym}
\allowdisplaybreaks[3]

\lstset{
language=Python,
basicstyle=\footnotesize,
keywordstyle=\color{blue},
stringstyle=\color{red},
commentstyle=\color{green!50!black},
showstringspaces=false,
breaklines=true,
frame=single,
numbers=left,
numberstyle=\tiny\color{gray},
stepnumber=1,
keepspaces=true,
extendedchars=true,
inputencoding=utf8
}

\newcommand{\num}[1]{#1}

\begin{document}
\sloppy
\thispagestyle{empty}
\centering{
Министерство образования и науки Российской Федерации\\
Федеральное государственное автономное образовательное учреждение высшего образования\\
«Санкт-Петербургский политехнический университет Петра Великого»\\[10pt]
Институт компьютерных наук и кибербезопасности\\
Высшая школа технологий искусственного интеллекта\\
02.03.01 Математика и компьютерные науки\\
\vspace{\fill}
{\large
Отчёт по дисциплине\\
{\Large \bfseries Вычислительная математика}\\
Лабораторная работа №4\\
{\Large <<Дискретное преобразование Фурье>>}\\
Вариант 18\\
\vspace{\fill}
}
{
\flushright
{\bfseries Выполнил:}\\
студент группы 5130201/40003\\
Джабраилов А.В.\\
\hspace{9cm}Подпись: \hrulefill \\
{\bfseries Принял:}\\
К. физ.-мат. наук, доц., доц.\\
Пак В.Г.\\
\hspace{9cm}Подпись: \hrulefill \\
\hspace{9cm}<<\rule{1.5cm}{0.4pt}>> \hrulefill\ 2026~г.
}
\centering{
Санкт-Петербург, 2026
}
}
\justifying
\newpage
\sloppy



\end{document}